\documentclass[paper=a4,fontsize=10pt,DIV11,BCOR10mm]{scrartcl}


\usepackage[english]{babel}
\usepackage[utf8]{inputenc}
\usepackage[T1]{fontenc}
\usepackage{lmodern}

\usepackage{amssymb}
\usepackage{amsmath}
\usepackage{amsthm}
\usepackage{cancel}

\usepackage{graphicx}

\usepackage{url}
\usepackage{eurosym}


\usepackage{url}



\newcommand{\abs}[1]{\left\lvert#1\right\rvert}
\newcommand{\norm}[1]{\left\lVert#1\right\rVert}
\DeclareMathOperator{\tr}{tr}




\titlehead{Technische Universität Berlin -- Fachgebiet Maschinelles Lernen\hfill \parbox[t]{2cm}{\includegraphics[width=2cm]{../TU_Logo_kurz_RGB_rot}}}

\begin{document}

\title{Maschinelles Lernen 1 - Assignment 5\\
\small{Technische Universität Berlin}}


\author{\small{Christoph Conrads (315565)}\and \small{Antje Relitz (327289)}  \and \small{Benjamin Pietrowicz (332542)} \and \small{Mitja Richter (324680)} }

\date{WS 2013/2014}

\maketitle


\section{Probabilistic PCA}
\subsection*{a)}
We want to show that $x \sim \mathcal{N}(\mu,C)$. In the proof we use the fact that the characteristic function of a multivariate Gaussian distribution $y \sim \mathcal{N}(\mu_0,\Sigma_0)$ is given by 
\begin{equation*}
\phi_y(t)=e^{-\frac{1}{2}i t^T \mu_0 + t^T \Sigma_0 t}
\end{equation*}
as well as the fact that the characteristic function of a sum of random variables is the product of characteristic functions of those random variables.\\

One more precondition that we will need is the fact that if $y \sim \mathcal{N}(\mu_0,\Sigma_0)$, then for a matrix $W$, the distribution of $Wy$ is given by $Wy \sim \mathcal{N}(W\mu_0,W\Sigma_0W^T)$.\\

To simplify notation, we will write, that $\mu$ is Gaussian distributed with mean $\mu$ and covariance matrix \underline{\underline{$0$}}. (This would theoretically cause problems, because the covariance matrix is supposed to be invertible, but we could also set the covariance matrix to be $\lambda I$ and then look at the limit $\lim\limits_{\lambda \rightarrow 0}\lambda I$, which would not change the result).\\

Therefore the individual characteristic functions are given by
\begin{align*}
\phi_{Wh}(t) &= \exp\{{-\frac{1}{2}i t^T W\underline{0} + t^T W I W^T t}\}\\
\phi_\mu(t) &= \exp\{{-\frac{1}{2}i t^T \underline{0} + t^T \underline{\underline{0}} t}\}\\
\phi_\epsilon(t) &= \exp \{{-\frac{1}{2}i t^T \underline{0} + t^T \sigma^2 I t}\}\text{.}\\
\end{align*}

With this preliminaries we do now obtain:
\begin{align*}
\phi_x(t) &= \phi_{Wh}(t) \cdot \phi_\mu(t) \cdot \phi_\epsilon(t)\\
&= \exp\{{-\frac{1}{2}i t^T W\underline{0} + t^T W I W^T t}\} \cdot \exp\{{-\frac{1}{2}i t^T \mu + t^T \underline{\underline{0}} t}\} \cdot \exp \{{-\frac{1}{2}i t^T \underline{0} + t^T \sigma^2 I t}\}\\
&=\exp\{{-\frac{1}{2}i t^T \mu + t^T W W^T t + t^T \sigma^2 I t}\}\\
&=\exp\{{-\frac{1}{2}i t^T \mu + t^T( W W^T +\sigma^2 I)t}\}\text{.}\\
\end{align*}
Therefore $x$ is multivariate Gaussian distributed with mean $\mu$ and covariance matrix $C=W W^T +\sigma^2 I$.

\subsection*{b)}
The covariance matrix is given by $C=W W^T +\sigma^2 I$. It can easily be seen that $C$ is symmetric and positive semi-definite (for any choice of $\sigma$). As $\sigma$ increases, the matrix $C$ becomes diagonally dominant and the diagonal entries become positive (since we add positive terms to the diagonal of $C$ only). We can now apply a result from linear algebra, that tells us that a symmetric diagonally dominant matrix with positive diagonal is also positive definite. So if the parameter $\sigma$ is large enough, the covariance matrix $C$ is positive definite and therefore invertible.\\

Since $WW^T$ is symmetric, there exist $d$ distinct eigenvectors $v_1,\dots, v_d$ of $WW^T$ corresponding to the eigenvalues $\mu_1,\dots \mu_d$ of $WW^T$ respectively, where the eigenvalues are ordered, such that $\mu_1 \geq \mu_2 \geq \dots \geq \mu_d$. 
Since 
\begin{equation*}
Cv_i=(WW^Tv_i+\sigma^2 I v_i)= \mu_i v_i +\sigma^2 v_i = (\mu_i+\sigma^2)v_i
\end{equation*}
holds, every eigenvector of $WW^T$ is also an eigenvector of $C$. Since $C$ is of dimension $d \times d$, $C$ can not have more than $d$ eigenvectors and therefore each eigenvector of $C$ is also an eigenvector of $WW^T$. We see that $\sigma^2$ is an additive constant, that is added to each eigenvalue $\mu_i$ and therefore the eigenvector $v_k$ corresponding to $\mu_k+\sigma^2$ ($k$-th principal component of $C$) is the same as the eigenvector of $WW^T$ corresponding to $\mu_k$, which is the $k$-th principal component of $WW^T$. Therefore the parameter $\sigma^2$ does not change the ordering of the principal components.

\subsection*{c)}
The complexity of $WW^T$ (and therefore the complexity of $C$ as well) increases with increasing $q$. This is due to the fact that we can interprete $WW^T$ as a multiple of the covariance matrix of measured datapoints $[w_1,\dots, w_q]=W$. As $q$ increases, the number of data points increases and the complexity of the covariance matrix $WW^T$ increases.\\

The idea behind probabilistic PCA is that the $d$-dimensional observed data is actually linearly dependent on some $q$-dimensional random distribution (and shifted by $\mu$ and falsified by some noise). Therefore it does not make sense to choose $q$ bigger or equal than $d$, since then we would try to explain our $d$-dimensional observation by a multivariate random distribution of dimension $d$ or more. This would lead to overfitting and therefore $q$ should be chosen to be smaller than $d$.




\section{Generalized Rayleigh Quotient}

\subsection*{a)}

Let $S_W, S_B \in \mathbb{R}^{d,d}$ where $S_W$ is symmetric positive definite. Let $A \in \mathbb{R}^{d,d}$ so that $A^T A = S_W$. Then
\[ J(w) = \frac{w^T S_B w}{w^T S_W w} = \frac{w^T S_B w}{(w^T A^T) (A w)} = \frac{w^T S_B w}{(A w)^T (A w)} \]
Let $v := A w$, then
\[ \bar{J}(v) = \frac{w^T A^T A^{-T} S_B A^{-1} A w}{v^T v} = \frac{v^T A^{-T} S_B A^{-1} v}{v^T v} = \frac{v^T M v}{v^T v} \]
where $M := A^{-T} S_B A^{-1}$.



\subsection*{b)}

We assume that there is a maximum. Let $w \neq 0$ be a maximizer of $\bar{J}$ and let $v = w / \norm{w} \Rightarrow \norm{v} = 1$. Thus
\[ \max_u \bar{J}(u) = \bar{J}(w) = \frac{w^T M w}{w^T w} = \frac{\norm{w}^2 v^T M v}{\norm{w}^2 v^T v} = \frac{v^T M v}{v^T v} = \bar{J}(v) \]
Therefore $v$ is a maximizer as well.




\subsection*{c)}

Using b) we look for vectors $v$ that maximize $\bar{J}$ subject to $\norm{v}_2 = 1$. Because the scaling of the vector has no influence on the function value $\bar{J}$ we choose $f(x) = x^T M x$ and $g(x) = x^T x - 1$. We are looking for
\[ \max_v f(v) - \lambda g(v) \text{ subject to } g(v) = 0 \]
The derivatives give
\begin{flalign*}
	&\frac{\partial f}{\partial v} = 2 M v \\
	&\frac{\partial g}{\partial v} = 2 v \\
\end{flalign*}
A necessary condition for the maximum is therefore
\[ 2 M v - 2 \lambda v = 0 \Leftrightarrow M v = \lambda v \]
which is maximal for the largest eigenvalue of $M$.



\subsection*{d)}

The function $J$ computes the eigenvalue $\lambda$ belonging to the equation
\[ S_B w = \lambda S_W w \]
If we multiply both sides with $S_W^{-1}$ we see that eigenvectors and eigenvalues remain unchanged. Therefore the largest eigenvalue of this problem
\[ S_W^{-1} S_B w = \lambda w \]
is also the maximizer of the generalized Rayleigh Quotient function $J(w)$.

\section{Extending Fisher Linear Discriminant}

\subsection*{a)}
Since the class means lie equally spaced on a one-dimensional line with most within-class variance othogonal to that line, there are two possible influences in this example. First of all we have three elliptic planes that are connected only through their middlepoint by that line. The line is orthogonal to these elliptic planes. The size of the ellipctic planes may vary as well as the distance between those planes with the restriction that the outer planes have the same distance to the inner plane. \\

The expected behavior in the extension is as follows:\\
We assume that our data is centralized with respect to the inner plane and the one-dimensional line is the x-axis itself. Therefore the distance to the other planes is $d$. With respect to our coordinate system the centerpoint of the first plane has the x-coordinate $-d$, the centerpoint of the second plane has the x-coordinate $0$ and the centerpoint of the third plane has the x-coordinate $d$. When we take a look at the between-class variance \mbox{$S_B = \sum_{i=1}^C\sum_{j=1}^C\left\Vert m_i - m_j \right\Vert^2 \cdot 1_{i>j}$} and we know that only the x-coordinates differ, we ultimatively know that: \[S_B = \left\Vert -d - 0 \right\Vert^2 + \left\Vert -d - d \right\Vert^2 + \left\Vert 0 - d\right\Vert^2 = d^2 + 4d^2 + d^2 = 6d^2 .\] This means our between-class variance grows quadratically with the distance $d$.\\

The within-class variance is defined as \mbox{$S_W = \sum_{i=1}^C S_i$, where $S_i = \sum_{k=1}^n (x_{k,i} -m_i)(x_{k,i} - m_i)^T $}. This means we have a sum of a sum. Since we have well-balanced classes the within-class variance is in the same order of magnitude and therefore does not differ too much. So the behavior in the sum should be $S_W \approx C \cdot S_i$ for any given class $i$.

\subsection*{b)}
The only difference to the first example is that the elliptic planes are not orthogonal to the one-dimensional line but parallel. To be more precisely in this example the line is in all of the elliptic planes and therefore a part of it. Additionally it goes again through the respective middlepoint.\\

The expected behavior is the same as in the first example since the spatial orientation of each class does not matter for scalar values.

\subsection*{c)}
The class means are at the corners of an isosecles triangle with small distinct angle. This means again that the inner class means $m_2$, the one at corner with the big angle, has the distance $d$ to the other class means $m_1$ and $m_3$, because we have an isoscele triangle. While the outer classes, the ones at the small distinct angle, have the distance $h$ to each other. We define the distance from $m_2$ to the line $h$ as $a$. The between-class variance for the case between class 1 and 3 is the only unknown distance with respect to $d$. Having an isosceles triangle makes the median line of $h$ the distance from $h$ to $m_2$ our defined $a$. Additionally our median line $a$ separates the triangle into two rectangular triangles, which means we can use the Pythagorean theorem. We know that $ d^2 = (\frac{h}{2})^2 + a^2$. For $\lim_{a \to 0} \frac{h}{2} = d \rightarrow \lim_{a \to 0} h = 2d $ and we have a one-dimensional line with all of our class means on that line with the distance $2d$ between the class means $m_1$ and $m_3$. Therefore we can assume $\lim_{a \to 0}h \approx 2d$. We can make this assumption $a \to 0$ because we have small distinct angles. This leads to the same result as in example 1 $S_B \approx 6d^2$.
Again, we have well-balanced classes which also leads to the same result as in example 1, that $S_W \approx C \cdot S_i$ for any class $i$.

%\begin{thebibliograpy}
%	\bibitem{DHS2000} R. O. Duda, P. E. Hart and D. G. Stork. \emph{Pattern Classification}. 2nd ed. New York, NY, USA: Wiley-Interscience, 2000. ISBN: 0-4710-5669-3
%\end{thebibliograpy}

\end{document}
